\section{Manifesto / LLM Results}\label{sec:v8-manifesto-results}

The Manifesto application asks whether a learned C-Tree can preserve a
social-science measurement target while exposing readable local states. The
source documents are party election programs from the Manifesto Project
\citep{ManifestoProject2025a,ManifestoCorpus2025}. The benchmark organized by
\citet{BenoitEtAl2025} matches $235$ party platforms to expert-survey means on
six seven-point policy dimensions: economic policy, social policy,
immigration, European integration, environmental policy, and decentralization.
An expert-survey mean is the document-level target obtained by averaging expert
placements for a party platform on one of these dimensions. For a given
dimension, Pearson \(r\) is the correlation across held-out manifestos between
the C-Tree prediction and that dimension's expert-survey mean. Macro \(r\) is
the unweighted average of those six dimension-level Pearson correlations.

This evidence should be read as a root-observed corpus evaluation. The held-out
documents have full-document expert targets, and the tree root is evaluated
directly against those targets. That is already the root channel in the
C-TreePO objective. It is not, by itself, a completed node-level local-law
certificate; that stronger claim requires sampled \(f^\ast\) checks inside the
realized tree.

Each manifesto--dimension pair is one document and one policy dimension. A
held-out pair is reserved for final evaluation rather than training or model
selection, and it has a full-document expert target. The root channel in
Equation~\eqref{eq:v8-main-objective} is therefore directly observed on the
evaluation split. The local-law channel is trained through teacher traces and
proxy scores: a teacher trace is a recorded local supervision signal for a
leaf or merge, and \(\widehat f\) supplies local losses on leaves and merges.
The stronger \(f^\ast\)-node audit described in Section~\ref{sec:v8-objective}
remains the next measurement layer.

\subsection{Tree Setup and Splits}\label{sec:v8-benoit-tree-setup}

The C-Tree version keeps Benoit's document-level target and changes the
representation being scored. A flat baseline scores the manifesto text
directly. A tree run chunks the manifesto into contiguous leaves, maps each
leaf into a readable policy-evidence state, recursively merges adjacent
states, and scores the root state. The f/g ladder alternates the state map and
readout: with \(\widehat f\) fixed, train \(g\) against node targets; with
\(g\) fixed, train \(\widehat f\) on the summaries produced by the current
tree. In this shorthand, \(g\) is one learned summarization/state map used
through the tree, and \(f\) or \(\widehat f\) is the readout used to score
states; a ladder stage such as \(f^2g^1\) means the second readout update
paired with the first state-map update.

The current f/g artifacts use $218$ manifestos split into $140$ training
documents, $30$ reserved documents used for model selection and diagnostics,
and $48$ held-out test documents used only for final reporting. Reported
metrics are held-out Pearson correlations against expert-survey means. The
generated figures and tables live under \texttt{assets/benoit/}.

\subsection{Single-Dimension Ladders}\label{sec:v8-benoit-single-dim}

Single-dimension training uses one \((\widehat f,g)\) pair per policy
dimension. This is the per-dimension ceiling for the tree representation
before asking one shared summary to serve all six dimensions.

\begin{figure}[!htbp]
  \centering
  \includegraphics[width=\linewidth]{manifesto_singledim_per_dim_live.pdf}
  \caption{\textbf{Single-dimension f/g ladders, all six policy dimensions.}
  External Pearson \(r\) against expert-survey means on the $48$-document
  held-out split, plotted against leaf token size. External means the score is
  computed on documents not used for training or model selection. Reference
  lines are from
  \citet{BenoitEtAl2025}.}
  \label{fig:v8-benoit-singledim-per-dim}
\end{figure}

Figure~\ref{fig:v8-benoit-singledim-per-dim} shows that separate
dimension-specific trees recover the published frontier range on most
dimensions. Economic policy, EU integration, immigration, and social policy
clear the proprietary ensemble reference in the best cells. Decentralization
is the stress dimension: it improves under single-dimension training, but its
audit gap remains high. Here an audit gap is the local diagnostic difference
between the teacher-trace target and the proxy/readout behavior, so a high
value marks the teacher/proxy channel as the limiting piece.

\subsection{Universal Summary}\label{sec:v8-benoit-universal-summary}

The universal-summary run trains one shared \(g\) to produce a single summary
per manifesto. Six per-dimension scorers then read that shared state. This is
the LLM application closest to the C-TreePO objective: root expert anchors
preserve the document target, while teacher/proxy local-law losses push one
state map to retain evidence for multiple downstream readouts.

\begin{table}[t]
  \centering
  \resizebox{\linewidth}{!}{% Generated table; do not edit by hand.
% Metric: Pearson r (higher better)
\begin{tabular}{lrrrrrrrr}
\toprule
Method & Econ. & Social & Immig. & EU & Env. & Decent. & Macro & \#/6 \\
\midrule
\multicolumn{9}{l}{\textit{Benoit et al. (2025) reference}} \\
Proprietary ensemble, 18 scores (Fig 1) & $+0.870$ & $+0.920$ & $+0.890$ & $+0.910$ & $+0.820$ & $+0.490$ & $+0.817$ & 6/6 \\
Split-expert reference (Table 3) & $+0.880$ & $+0.910$ & $+0.880$ & $+0.950$ & $+0.840$ & $+0.780$ & $+0.873$ & 6/6 \\
LLaMA-3.3-70B (Table 6) & $+0.840$ & $+0.870$ & $+0.860$ & $+0.860$ & $+0.680$ & $+0.400$ & $+0.752$ & 6/6 \\
DeepSeek-V3 (Table 6) & $+0.840$ & $+0.870$ & $+0.890$ & $+0.860$ & $+0.790$ & $+0.450$ & $+0.783$ & 6/6 \\
Gemma-3-27B-IT (Table 6) & $+0.860$ & $+0.860$ & $+0.890$ & $+0.840$ & $+0.860$ & $+0.450$ & $+0.793$ & 6/6 \\
\midrule
\multicolumn{9}{l}{\textit{Ours: per-dim pipeline $\times$ leaf size (Gemma-4-31B-NVFP4, 1 summarizer + 1 scorer per dim)}} \\
leaf = 64 K chars ($\approx$16K tokens) & $+0.916$ & $+0.817$ & $+0.888$ & $+0.924$ & $+0.854$ & $+0.489$ & $+0.815$ & 6/6 \\
leaf = 32 K chars ($\approx$8K tokens) & $+0.929$ & $+0.857$ & $+0.886$ & $+0.916$ & $+0.877$ & $+0.480$ & $+0.824$ & 6/6 \\
leaf = 24 K chars ($\approx$6K tokens) --- full test n$\approx$215 & $+0.892$ & $+0.840$ & $+0.867$ & $+0.896$ & $+0.814$ & $+0.464$ & $+0.796$ & 6/6 \\
leaf = 16 K chars ($\approx$4K tokens) & $+0.925$ & $+0.847$ & $+0.898$ & $+0.909$ & $+0.831$ & $+0.537$ & $+0.824$ & 6/6 \\
leaf =  8 K chars ($\approx$2K tokens) & $+0.939$ & $+0.869$ & $+0.884$ & $+0.901$ & $+0.839$ & $+0.543$ & $+0.829$ & 6/6 \\
\midrule
\multicolumn{9}{l}{\textit{Ours: combined pipeline $\times$ leaf size (one shared summarizer with union rubric, 6 scores)}} \\
leaf = 64 K chars & $+0.910$ & $+0.867$ & $+0.897$ & $+0.910$ & $+0.795$ & $+0.458$ & $+0.806$ & 6/6 \\
leaf = 32 K chars & $+0.917$ & $+0.849$ & $+0.908$ & $+0.889$ & $+0.749$ & $+0.438$ & $+0.792$ & 6/6 \\
leaf = 24 K chars (full test n=229) & $+0.868$ & $+0.850$ & $+0.880$ & $+0.902$ & $+0.809$ & $+0.396$ & $+0.784$ & 6/6 \\
leaf = 16 K chars & $+0.919$ & $+0.851$ & $+0.916$ & $+0.889$ & $+0.808$ & $+0.441$ & $+0.804$ & 6/6 \\
leaf =  8 K chars & $+0.927$ & $+0.874$ & $+0.911$ & $+0.917$ & $+0.801$ & $+0.413$ & $+0.807$ & 6/6 \\
\midrule
\multicolumn{9}{l}{\textit{Ours: tiny-leaf extensions of the chunk sweep (Gemma-4)}} \\
tree, leaf = 4 K chars & $+0.932$ & $+0.886$ & $+0.885$ & $+0.931$ & $+0.838$ & $+0.580$ & $+0.842$ & 6/6 \\
tree, leaf = 2 K chars & $+0.924$ & $+0.867$ & $+0.887$ & $+0.906$ & $+0.793$ & $+0.584$ & $+0.827$ & 6/6 \\
tree, leaf = 1 K chars (stress test, depth $\geq$6) & $+0.933$ & $+0.859$ & $+0.914$ & $+0.925$ & $+0.838$ & $+0.493$ & $+0.827$ & 6/6 \\
\midrule
\multicolumn{9}{l}{\textit{Ours: concat-no-merge (chunks summarized independently, joined, scored --- tests whether the merge step carries signal)}} \\
concat, leaf = 32 K chars & $+0.929$ & $+0.856$ & $+0.897$ & $+0.920$ & $+0.848$ & $+0.571$ & $+0.837$ & 6/6 \\
concat, leaf = 16 K chars (default) & $+0.931$ & $+0.823$ & $+0.909$ & $+0.912$ & $+0.866$ & $+0.523$ & $+0.827$ & 6/6 \\
concat, leaf =  8 K chars & $+0.932$ & $+0.849$ & $+0.882$ & $+0.915$ & $+0.850$ & $+0.594$ & $+0.837$ & 6/6 \\
\midrule
\multicolumn{9}{l}{\textit{Ours: flat baseline (no chunk, no summary; truncate text and score) --- tests whether tree is needed at all}} \\
flat, truncation = 48 K chars & $+0.936$ & $+0.821$ & $+0.926$ & $+0.889$ & $+0.806$ & $+0.587$ & $+0.827$ & 6/6 \\
flat, truncation = 24 K chars (default) & $+0.929$ & $+0.847$ & $+0.889$ & $+0.882$ & $+0.829$ & $+0.542$ & $+0.820$ & 6/6 \\
flat, truncation = 12 K chars & $+0.929$ & $+0.844$ & $+0.938$ & $+0.880$ & $+0.779$ & $+0.564$ & $+0.822$ & 6/6 \\
flat, truncation =  6 K chars & $+0.926$ & $+0.811$ & $+0.967$ & $+0.881$ & $+0.886$ & $+0.243$ & $+0.786$ & 6/6 \\
\midrule
\multicolumn{9}{l}{\textit{Ours: scorer ablations (Gemma-4, Benoit's GPT-4o summaries held fixed)}} \\
1. Zero-shot scorer & $+0.832$ & $+0.886$ & $+0.858$ & $+0.905$ & $+0.668$ & $+0.311$ & $+0.743$ & 6/6 \\
2. Few-shot optimized scorer & $+0.822$ & $+0.892$ & $+0.853$ & $+0.912$ & $+0.627$ & $+0.297$ & $+0.734$ & 6/6 \\
3. Joint scorer baseline (shared across 6 dims) & $+0.837$ & $+0.881$ & $+0.852$ & $+0.904$ & $+0.658$ & $+0.314$ & $+0.741$ & 6/6 \\
4. Joint scorer, few-shot optimized & $+0.817$ & $+0.891$ & $+0.848$ & $+0.908$ & $+0.700$ & $+0.313$ & $+0.746$ & 6/6 \\
5. Joint scorer, reflective prompt-optimized & $+0.832$ & $+0.882$ & $+0.848$ & $+0.879$ & $+0.688$ & $+0.344$ & $+0.746$ & 6/6 \\
\midrule
\multicolumn{9}{l}{\textit{Ours: exact-model replication (Benoit's Gemma-3-27B-IT BF16)}} \\
Gemma-3-27B scorer-only, OUR scoring rubric & $+0.826$ & $+0.864$ & $+0.846$ & $+0.902$ & $+0.573$ & $+0.329$ & $+0.723$ & 6/6 \\
Gemma-3-27B scorer-only, exact Benoit rubric & $+0.814$ & $+0.895$ & $+0.852$ & $+0.852$ & $+0.616$ & $+0.359$ & $+0.731$ & 6/6 \\
Gemma-3-27B, raw-prompt Benoit format (bare-integer response) & $+0.782$ & $+0.901$ & $+0.854$ & $+0.874$ & $+0.590$ & $+0.368$ & $+0.728$ & 6/6 \\
\bottomrule
\end{tabular}
}
  \caption{\textbf{Pearson \(r\) against Benoit expert-survey means.}
  Reference rows reproduce the Benoit benchmark; C-Tree rows use the learned
  manifesto compression pipeline.}
  \label{tab:v8-benoit-headline}
\end{table}

Table~\ref{tab:v8-benoit-headline} gives the six-dimension benchmark. The
per-dimension tree reaches macro \(r=0.829\) at the reported character-leaf
setting, meaning the leaf size is controlled in characters rather than tokens.
The universal-summary tree reaches macro \(r=0.807\), with most of the gap
concentrated in decentralization. The result is consistent with the
state-sufficiency view: one summary can carry the shared evidence for several
dimensions, while dimensions with different evidence geometry need either more
capacity, different law weights, or their own state map.

\begin{figure}[!htbp]
  \centering
  \includegraphics[width=\linewidth]{manifesto_fg_combined_ladder_f1g0_f1g1.pdf}
  \caption{\textbf{Universal-summary run, per-dimension trajectories.}
  External Pearson \(r\) against expert-survey means on the $48$-document
  held-out split from one shared \((\widehat f,g)\) training trajectory.}
  \label{fig:v8-benoit-combined-ladder}
\end{figure}

\begin{table}[t]
  \centering
  \begin{tabular}{lrrrr}
    \toprule
    Dimension & Best \(r\) & Leaf tokens & Stage & \(n\) \\
    \midrule
    Economic & \(0.886\) & \(4{,}096\) & \(f^2g^1\) & \(47\) \\
    Social & \(0.886\) & \(1{,}024\) & \(f^2g^2\) & \(43\) \\
    Immigration & \(0.938\) & \(4{,}096\) & \(f^3g^3\) & \(29\) \\
    European integration & \(0.965\) & \(2{,}048\) & \(f^1g^1\) & \(35\) \\
    Environment & \(0.795\) & \(1{,}024\) & \(f^1g^1\) & \(42\) \\
    Decentralization & \(0.461\) & \(8{,}096\) & \(f^1g^1\) & \(42\) \\
    \bottomrule
  \end{tabular}
  \caption{\textbf{Best per-dimension cells under the universal summary.}
  Leaf tokens is the token budget for one leaf, stage is the f/g ladder state,
  and \(n\) is the number of held-out documents with usable labels for that
  dimension. A common deployed row would choose one leaf size and one stage;
  this table shows the best cell per dimension to diagnose state-sharing
  limits.}
  \label{tab:v8-benoit-joint-dim-best}
\end{table}

Figure~\ref{fig:v8-benoit-combined-ladder} and
Table~\ref{tab:v8-benoit-joint-dim-best} show the same pattern at the
trajectory level. Economic, social, immigration, EU integration, and
environmental policy remain close to their single-dimension ceilings.
Decentralization loses more under shared compression, which points to
multi-task interference rather than a generic failure of tree compression.

\subsection{Local Diagnostics}\label{sec:v8-benoit-local-diagnostics}

The available local diagnostics are teacher-trace residuals, not a completed
human node audit. A teacher-trace residual is the discrepancy between a
recorded teacher local target and the current proxy/readout behavior at the
corresponding node. These diagnostics still serve the v8 objective story: the
proxy \(\widehat f\) produces local-law losses everywhere, while future sampled
\(f^\ast\) node calls can correct the proxy channel using
Equation~\eqref{eq:v8-corrected-node-loss}.

\begin{figure}[!htbp]
  \centering
  \includegraphics[width=\linewidth]{manifesto_fg_combined_audit_gap.pdf}
  \caption{\textbf{Universal-summary local diagnostic gaps.}
  Per-dimension teacher-trace gaps from the universal-summary run.}
  \label{fig:v8-benoit-combined-audit-gap}
\end{figure}

\begin{figure}[!htbp]
  \centering
  \includegraphics[width=\linewidth]{manifesto_singledim_per_dim_live_audit_gap.pdf}
  \caption{\textbf{Single-dimension local diagnostic gaps.}
  Teacher-trace gap for the single-dimension ladders in
  Figure~\ref{fig:v8-benoit-singledim-per-dim}.}
  \label{fig:v8-benoit-single-audit-gap}
\end{figure}

The local diagnostic figures identify where proxy supervision is least aligned
with the held-out expert target. Under the v8 objective, that is exactly where
sampled \(f^\ast\) node labels have the highest value: they estimate and
correct \(b_v=\widehat\ell_v-\ell_v^\ast\) rather than merely adding another
root score.

\subsection{Takeaway}\label{sec:v8-benoit-takeaway}

The Manifesto results instantiate the stack from
Section~\ref{sec:v8-objective}. Expert-survey means supply the observed root
channel, teacher traces and proxy readouts supply dense local training geometry
for \(g\), and the diagnostic gaps identify where sampled \(f^\ast\) node
calls would most improve the corrected local-law channel. The current evidence
therefore supports the root-observed LLM compression claim and sets up the
stronger node-level audit, rather than treating the proxy diagnostics as a
separate objective.
